% ===== ch12 : 原文 main.tex 第 4216--4655 行（6c 扭 Künneth 公式）=====
% 本文件是原文 LaTeX 源码的逐行分片（去掉整段注释掉的旧稿），仅供翻译时对照；最终译文写到 src/chunks/ch12.tex。
\subsection{Twisted Künneth formula.}
We end this section with the K\"unneth theorem. Suppose that $X$ and $Y$ are $G$- and $H$-schemes that admits an open cover $\mathfrak{U},\mathfrak{V}$ of $G$-invariant, resp. $H$-invariant open affine subschemes. Recall that $X\times Y$ admits an open cover by affine schemes of the form $U\times V\in\mathfrak{U}\times\mathfrak{V}$, which is clearly $G\times H$-invariant, 
Then we can define the equivariant K\"unneth functor
$$
% ---- 原文行 4220 ----
\Perf_G(X;\mathfrak{U})_{\dg}
\otimes
\Perf_H(Y;\mathfrak{V})_{\dg}
\xrightarrow{\Ku}
\Perf_{G\times H}(X\times Y;\mathfrak{U}\times\mathfrak{V})_{\dg}
$$
on dg enhancements as follows:
\begin{defn}
$\Ku$ is a functor such that
\begin{itemize}
\item
Objects:
It maps an object
$
\big((\mathcal{F}^{\bullet},\theta^{\bullet}),(\mathcal{G}^{\bullet},\eta^{\bullet})\big)
$
%in
%$
%\Perf_G(X;\mathfrak{U})_{\dg}
%\otimes
% ---- 原文行 4240 ----
%\Perf_H(Y;\mathfrak{V})_{\dg}
%$
to an object
%in
%$
%\Perf_{G\times H}(X\times %Y;\mathfrak{U}\times\mathfrak{V})_{\dg}
%$
%given by
%$$
%\Big(
%\bigoplus_{n\in\mathbb{Z}}
%\big(
%(\mathcal{F}^{\bullet},\theta^{\bullet})
%\boxtimes
%(\mathcal{G}^{\bullet},\eta^{\bullet})
%\big)^n,
%\
%d_{\mathcal{F}^{\bullet}\boxtimes\mathcal{G}^{\bullet}}
%:=
%\big(
% ---- 原文行 4260 ----
%d_{\mathcal{F}^{\bullet}\boxtimes\mathcal{G}^{\bullet}}^n
%:
%(\mathcal{F}^{\bullet}\boxtimes\mathcal{G}^{\bullet})^n
%\rightarrow
%(\mathcal{F}^{\bullet}\boxtimes\mathcal{G}^{\bullet})^{n+1}
%\big)_{n\in\mathbb{Z}}
%\Big)
%$$
$$
\Big(
\bigoplus_{n\in\mathbb{Z}}
\big(
(\mathcal{F}^{\bullet}\boxtimes\mathcal{G}^{\bullet})^n,
(\theta^{\bullet}\boxtimes\eta^{\bullet})^n
\big),
\
d_{\mathcal{F}^{\bullet}\boxtimes\mathcal{G}^{\bullet}}
:=
\big(
d_{\mathcal{F}^{\bullet}\boxtimes\mathcal{G}^{\bullet}}^n
% ---- 原文行 4280 ----
:
(\mathcal{F}^{\bullet}\boxtimes\mathcal{G}^{\bullet})^n
\rightarrow
(\mathcal{F}^{\bullet}\boxtimes\mathcal{G}^{\bullet})^{n+1}
\big)_{n\in\mathbb{Z}}
\Big)
$$
%$
%\Big(
%\mathcal{F}^{\bullet}\boxtimes\mathcal{G}^{\bullet}
%:=
%\bigoplus_{n\in\mathbb{Z}}(\mathcal{F}^{\bullet}\boxtimes\mathcal{G}^{\bullet})^n,
%\
%\theta^{\bullet}\boxtimes\eta^{\bullet}
%:=
%\bigoplus_{n\in\mathbb{Z}}(\theta^{\bullet}\boxtimes\eta^{\bullet})^n
%\Big)
%$
(see Definition \ref{dfn:GEquivQCohdg}),
where
% ---- 原文行 4300 ----
$$
\big(
(\mathcal{F}^{\bullet}\boxtimes\mathcal{G}^{\bullet})^n,
(\theta^{\bullet}\boxtimes\eta^{\bullet})^n
\big)
:=
\bigoplus_{i+j=n}
(
\mathcal{F}^i\boxtimes\mathcal{G}^j,
\
\theta^i\boxtimes\eta^j
)
\quad\text{for each $n\in\mathbb{Z}$}
$$
(note that
$
(\mathcal{F}^i,\theta^i)
\boxtimes
(\mathcal{G}^j,\eta^j)
=
% ---- 原文行 4320 ----
(
\mathcal{F}^i\boxtimes\mathcal{G}^j,
\
\theta^i\boxtimes\eta^j
)
\in\Mod_{G\times H}(X\times Y)
$
was defined in
Definition \ref{defn:ExternalTensorProductOfEquivariantModules}),
and
the differential is determined on each summand
$\mathcal{F}^i\boxtimes\mathcal{G}^j$
(where $i+j=n$)
by
\[
d_{\mathcal F^\bullet\boxtimes\mathcal G^\bullet}^n
\big|_{\mathcal F^i\boxtimes\mathcal G^j}
:=
d_{\mathcal F}^i\boxtimes\id_{\mathcal G^j}
+
% ---- 原文行 4340 ----
(-1)^i\id_{\mathcal F^i}\boxtimes d_{\mathcal G}^j
:
\mathcal{F}^i\boxtimes\mathcal{G}^j
\rightarrow
(\mathcal{F}^{i+1}\boxtimes\mathcal{G}^j)
\oplus
(\mathcal{F}^i\boxtimes\mathcal{G}^{j+1}).
\]

\item
Morphisms:
%For homogeneous morphisms
%\begin{align*}
%f
%&\in
%\Hom_{\VB_G(U_{i_0\cdots i_p})}
%\Big(
%(\mathcal F^{\objectdegree},\theta^{\objectdegree})
%  |_{U_{i_0\cdots i_p}},
%(\mathcal F'^{\objectdegree+q},\theta'^{\objectdegree+q})
% ---- 原文行 4360 ----
%  |_{U_{i_0\cdots i_p}}
%\Big)
%\\
%&\subseteq
%\hom_{\Perf_G(X;\mathfrak U)_{\dg}}^k
%\big(
%(\mathcal F^\bullet,\theta^\bullet),
%(\mathcal F'^\bullet,\theta'^\bullet)
%\big)
%\quad
%(k=p+q)
%\\
%\text{and}\quad
%g
%&\in
%\Hom_{\VB_H(V_{j_0\cdots j_r})}
%\Big(
%(\mathcal G^{j},\eta^{j})
%  |_{V_{j_0\cdots j_r}},
%(\mathcal G'^{j+s},\eta'^{j+s})
% ---- 原文行 4380 ----
%  |_{V_{j_0\cdots j_r}}
%\Big)
%\\
%&\subseteq
%\hom_{\Perf_H(Y;\mathfrak V)_{\dg}}^l
%\big(
%(\mathcal G^\bullet,\eta^\bullet),
%(\mathcal G'^\bullet,\eta'^\bullet)
%\big)
%\quad
%(l=r+s),
%\end{align*}
%their tensor product $f\otimes g$
%is mapped to
%\begin{align*}
%f\boxtimes g
%&\in
%\Hom_{\VB_{G\times H}(U_{i_0\cdots i_p}\times V_{j_0\cdots j_r})}
%\Big(
%(\mathcal{F}^i\boxtimes\mathcal{G}^j,
% ---- 原文行 4400 ----
%\theta^i\boxtimes\eta^j)|
%_{U_{i_0\cdots i_p}\times V_{j_0\cdots j_r}},
%(\mathcal{F}'^{i+q}\boxtimes\mathcal{G}'^{j+s},
%\theta'^{i+q}\boxtimes\eta'^{j+s})|
%_{U_{i_0\cdots i_p}\times V_{j_0\cdots j_r}}
%\Big)
%\\
%&\subseteq
%\hom_{\Perf_{G\times H}
%(X\times Y;\mathfrak U\times\mathfrak V)_{\dg}}^{k+l}
%\big(
%(\mathcal F^\bullet,\theta^\bullet)
% \boxtimes(\mathcal G^\bullet,\eta^\bullet),
%(\mathcal F'^\bullet,\theta'^\bullet)
% \boxtimes(\mathcal G'^\bullet,\eta'^\bullet)
%\big)
%\end{align*}
%(recall Definition \ref{defn:ExternalTensorProductOfEquivariantModules}).
For homogeneous morphisms $f\in C_{\otimes}^{p+q}$ and $g\in C_{\otimes}^{r+s}$,
% \adjustbox{width=160mm}{
% ---- 原文行 4420 ----
% $
% f
% :=
% \Big(
% f_{i_0\cdots i_p}^i
% \in
% \Hom_{\VB_G(U_{i_0\cdots i_p})}\big((\mathcal{F}^i,\theta^{\objectdegree})|_{U_{i_0\cdots i_p}},(\mathcal{F}'^{i+q},\theta'^{i+q})|_{U_{i_0\cdots i_p}}\big)
% \Big)_{
% \substack{
% i\in\mathbb{Z}
% \\
% i_0<\cdots<i_p
% }
% }
% \in
% \check{C}_G^p\big(
% \mathfrak{U},\mathcal{H}om_{\mathcal{O}_X}^q(
% (\mathcal{F}^{\bullet},\theta^{\bullet}),
% (\mathcal{F}'^{\bullet},\theta'^{\bullet})
% )
% ---- 原文行 4440 ----
% \big),
% $
% }
% \\
% \adjustbox{width=160mm}{
% $
% g
% :=
% \Big(
% g_{j_0\cdots j_r}^j
% \in
% \Hom_{\VB_H(V_{j_0\cdots j_r})}\big((\mathcal{G}^j,\eta^j)|_{V_{j_0\cdots j_r}},(\mathcal{G}'^{j+s},\eta'^{j+s})|_{V_{j_0\cdots j_r}}\big)
% \Big)_{
% \substack{
% j\in\mathbb{Z}
% \\
% j_0<\cdots<j_r
% }
% }
% \in
% ---- 原文行 4460 ----
% \check{C}_H^r\big(
% \mathfrak{V},\mathcal{H}om_{\mathcal{O}_Y}^s(
% (\mathcal{G}^{\bullet},\eta^{\bullet}),
% (\mathcal{G}'^{\bullet},\eta'^{\bullet})
% )
% \big),
% $
% }
their tensor product $f\otimes g$ is mapped to $T^{\ast} (f\otimes g)\in C_{\times}^{p+q+r+s}$, where $T^{\ast}$ is the map defined in section \ref{subsec:Eilenberg-Zilber-contraction-Cech}.
% \adjustbox{scale=0.75}{
% $
% f\boxtimes g
% :=
% \Big(
% (f\boxtimes g)_{\alpha_0\cdots \alpha_t}^k
% \in
% \displaystyle
% \bigoplus_{i+j=k}
% \Hom_{\VB_{G\times H}(W_{\alpha_0\cdots \alpha_t})}
% \big((\mathcal{F}^i\boxtimes\mathcal{G}^j,\theta^i\boxtimes\eta^j)|_{W_{\alpha_0\cdots \alpha_t}},(\mathcal{F}'^{i+q}\boxtimes\mathcal{G}'^{j+s},\theta'^{i+q}\boxtimes\eta'^{j+s})|_{W_{\alpha_0\cdots \alpha_t}}\big)
% ---- 原文行 4480 ----
% \Big)_{
% \substack{
% k\in\mathbb{Z}
% \\
% \alpha_0<\cdots<\alpha_t
% }
% }
% $}

% \adjustbox{scale=0.75}{
% $
% \phantom{f\boxtimes g:}
% \in
% \check{C}_{G\times H}^t
% \Big(
% \mathfrak{U}\times\mathfrak{V},\mathcal{H}om_{\mathcal{O}_{X\times Y}}^u
% \big(
% (\mathcal F^\bullet,\theta^\bullet)
%  \boxtimes(\mathcal G^\bullet,\eta^\bullet),
% (\mathcal F'^\bullet,\theta'^\bullet)
% ---- 原文行 4500 ----
%  \boxtimes(\mathcal G'^\bullet,\eta'^\bullet)
% \big)
% \Big)
% \qquad
% (t=p+r,\ u=q+s)
% $
% }

% (note that $p$, $q$, $r$, $s$, and hence $t$, $u$ are fixed),
% where
% $$
% \left.
% (f\boxtimes g)_{\alpha_0\cdots\alpha_t}^k
% \right|_{\mathcal{F}^i\boxtimes\mathcal{G}^j}
% :=
% \displaystyle
% %\sum_{i+j=k}
% (-1)^{qr+qs+si}
% \left.
% \left(
% ---- 原文行 4520 ----
% f_{i_0\cdots i_p}^i
% \boxtimes
% g_{j_p\cdots j_t}^j
% \right)
% \right|_{W_{\alpha_0\cdots\alpha_t}}
% $$
% \footnote{
% Here,
% we denote by
% $
% g_{j_{\sigma(0)}\cdots j_{\sigma(r)}}^j
% :=
% \operatorname{sign}(\sigma)
% g_{j_0\cdots j_r}^j
% $
% for each $j_0<\cdots <j_r$
% and
% $\sigma\in S_{r+1}$.
% A component with repeated indices is defined to be zero.
% The same convention applies to $f$.
% ---- 原文行 4540 ----
% }
% for every
% $\alpha_0=(i_0,j_0)<\cdots<\alpha_t=(i_t,j_t)$
% and
% $i$, $j\in\mathbb{Z}$ with $i+j=k$
% (recall Definition \ref{defn:ExternalTensorProductOfEquivariantModules}).
\item Higher structures:
we define the $n$-th K\"unneth functor 
\begin{align*}
\Ku^n\colon &\hom_{\Perf (X;\mathfrak{U})_{dg}\otimes\Perf(Y;\mathfrak{V})_{dg}}\left(\left(
    (\mathcal{F}_{n-1}^{\bullet},\theta_{n-1}^{\bullet}),(\mathcal{G}_{n-1}^{\bullet},\eta_{n-1}^{\bullet})\right),
    \left((\mathcal{F}_n^{\bullet},\theta_n^{\bullet}),(\mathcal{G}_{n}^{\bullet},\eta_n^{\bullet})\right)\right)
    \otimes\dotsb\\ 
    \otimes
    &\hom_{\Perf (X;\mathfrak{U})_{dg}\otimes\Perf(Y;\mathfrak{V})_{dg}}\left(\left(
    (\mathcal{F}_0^{\bullet},\theta_0^{\bullet}),(\mathcal{G}_0^{\bullet},\eta_0^{\bullet})\right),
    \left((\mathcal{F}_1^{\bullet},\theta_1^{\bullet}),(\mathcal{G}_1^{\bullet},\eta_1^{\bullet})\right)\right)\\
    \to &\hom_{\Perf_{G\times H} (X\times Y;\mathfrak{U}\times\mathfrak{V})_{dg}}
    \left((\mathcal{F}_0^{\bullet},\theta_0^{\bullet})\boxtimes (\mathcal{G}_0^{\bullet},\eta_0^{\bullet}),(\mathcal{F}_n^{\bullet},\theta_n^{\bullet})\boxtimes (\mathcal{G}_n^{\bullet},\eta_n^{\bullet})\right)
\end{align*}
% ---- 原文行 4560 ----
provided by the formula 
\begin{center}
\adjustbox{max width=\textwidth}{$\displaystyle
    \Ku^n (f_{n-1}\otimes g_{n-1},\dotsb, f_0\otimes g_0)=\sum_{\ell=1}^{n-1} (-1)^{(n-\ell -1)(t_{\ell}'+u_{\ell}'-1)}\mathbb{H}^{\ast}\left(
        \Ku^{n-\ell} (f_{n-1}\otimes g_{n-1},\dotsb, f_{\ell}\otimes g_{\ell})\circ\Ku^{\ell} (f_0\otimes g_0,\dotsb, f_{\ell-1}\otimes g_{\ell-1})
        \right),
$}\end{center}
where the operator $\mathbb{H}^{\ast}$ is defined in appendix \ref{sec:Eilenberg-Zilber-constraction}.
\end{itemize}
\end{defn}
We can check directly that
\begin{prp}
    The collection $(\Ku^n)_{n\in\mathbb{Z}_{\geq 0}}$ forms an $A_{\infty}$-functor from $\Perf_G (X;\mathfrak{U})_{dg}\otimes\Perf_H (Y;\mathfrak{V})_{dg}$ to $\Perf (X\times Y;\mathfrak{U}\times\mathfrak{V})_{dg}$.
\end{prp}
\begin{proof}
    The full verification will be lengthy, and we only present the verification for the first non-trivial relation to provide the pattern. From appendix \ref{sec:Eilenberg-Zilber-constraction} we see that $T^{\ast} S^{\ast} =\id$ and $S^{\ast} T^{\ast} -\id =\mathbb{H}^{\ast}d+d\mathbb{H}^{\ast}$, and hence 
    \begin{align*}
        d\Ku^2& (f_1\otimes g_1,f_0\otimes g_0)=d\mathbb{H}^{\ast} (\Ku (f_1\otimes g_1)\circ\Ku (f_0\otimes g_0))\\ 
        &=-\mathbb{H}^{\ast} (\Ku (d(f_1\otimes g_1))\circ \Ku (f_0\otimes g_0)) -(-1)^{\vert f_0\otimes g_0\vert }\mathbb{H}^{\ast} (\Ku(f_1\otimes g_1)\circ\Ku (d(f_0\otimes g_0)))\\
        &+T^{\ast} (S^{\ast}(\Ku (f_1\otimes g_1))\circ S^{\ast} (\Ku (f_0\otimes g_0))) -\Ku (f_1\otimes g_1)\circ\Ku (f_0\otimes g_0)\\ 
% ---- 原文行 4580 ----
        &=-(-1)^{\vert f_0\otimes g_0\vert }\Ku^2 (d(f_1\otimes g_1),f_0\otimes g_0) - \Ku^2(f_1\otimes g_1,d(f_0\otimes g_0))\\ 
        &+\Ku((f_1\otimes g_1)\circ (f_0\otimes g_0))-\Ku (f_1\otimes g_1)\circ\Ku (f_0\otimes g_0),
    \end{align*}
    which is exactly the $A_{\infty}$-relation for $\Ku^2$. 
\end{proof}
Now we are ready to prove the K\"unneth embedding.
\begin{thm}\label{thm:Kunneth-embedding}
The prescribed functor
$$
\Ku\colon\Perf_G(X;\mathfrak{U})_{\dg}
\otimes
\Perf_H(Y;\mathfrak{V})_{\dg}
\hookrightarrow
\Perf_{G\times H}(X\times Y;\mathfrak{U}\times\mathfrak{V})_{\dg}
$$

$$
(\mathcal{F},\mathcal{G})
\quad\mapsto\quad
\mathcal{F}\boxtimes\mathcal{G}
% ---- 原文行 4600 ----
~:=~
\pi_X^*\mathcal{F}
\underset{\mathcal{O}_{X\times Y}}{\otimes}
\pi_Y^*\mathcal{G}
$$
$$
f\otimes g
\quad\mapsto\quad
f\boxtimes g
~:=~
\pi_X^*f
\otimes
\pi_Y^*g
$$
is a quasi-embedding of dg-categories.
\end{thm}
\begin{proof}
    From Appendix \ref{sec:Eilenberg-Zilber-constraction} we know that $\Ku^1 =T^{\ast}$ is a quasi-equivalence of chain complexes with quasi-inverse $S^{\ast}$, and therefore after passing to homology category we know that $H^0 (\Ku)\colon H^0 (\Perf_G (X;\mathfrak{U})_{dg}\otimes\Perf_H (Y;\mathfrak{V})_{dg})\to H^0 (\Perf_{G\times H} (X\times Y;\mathfrak{U}\times\mathfrak{V})_{dg})$ is an embedding of categories. This proves that $\Ku$ is a quasi-embedding.
\end{proof}

% ---- 原文行 4620 ----
\begin{crl}
    $\Ku$ induces a quasi-embedding
$$
\Perf_G(X)_{\dg}
\underset{\widehat{G}}{\otimes}
\Perf_G(Y)_{\dg}
\hookrightarrow
\Perf_G(X\times Y)_{\dg}
$$

$$
(\mathcal{F},\mathcal{G})
\quad\mapsto\quad
\mathcal{F}\boxtimes\mathcal{G}
~:=~
\pi_X^*\mathcal{F}
\underset{\mathcal{O}_{X\times Y}}{\otimes}
\pi_Y^*\mathcal{G}
$$
$$
% ---- 原文行 4640 ----
f\otimes g
\quad\mapsto\quad
f\boxtimes g
~:=~
\pi_X^*f
\otimes
\pi_Y^*g.
$$
\end{crl}
\begin{proof}
From theorem \ref{thm:Kunneth-embedding} we get a quasi-embedding\\ 
$\Perf_G (X)_{dg}\otimes\Perf_G(Y)_{dg}\hookrightarrow\Perf_{G\times G} (X\times Y)_{dg}$. On the other hand, consider the diagonal embedding $\Delta\colon G\hookrightarrow G\times G$, its orthogonal subgroup action $\widehat{G}\cong (\Delta (G))^{\perp}\curvearrowright\Perf_{G\times G} (Y)$ preserves the subcategory $\Perf_G (X)_{dg}\otimes\Perf_G(Y)_{dg}$, and therefore from Corollary \ref{crl:RestrictionFunctionEm} and Proposition \ref{prp:twisted-tensor-equal-quotient} we get a quasi-embedding \begin{align*}
    \Perf_G (X)_{dg}\underset{\widehat{G}}{\otimes}\Perf_G(Y)_{dg}&\simeq\left(\Perf_G(X)\otimes\Perf_G(Y)_{dg}\right)\big/\widehat{G}\\ 
    &\hookrightarrow\Perf_{G\times G} (X\times Y)_{dg}\big/\widehat{G}\hookrightarrow\Perf_G (X\times Y)_{dg}.\qedhere
\end{align*}
\end{proof}
